\medskip

On considère la fonction $f$ définie sur l'intervalle [0~;~1] par 
\[f(x) = 2x\e^{-x}.\]

On admet que la fonction $f$ est dérivable sur l'intervalle [0~;~1]. 

\medskip

\begin{enumerate}
\item 
	\begin{enumerate}
		\item Résoudre sur l'intervalle [0~;~1] l'équation $f(x) = x$.
		\item Démontrer que, pour tout $x$ appartenant à l'intervalle [0~;~1], 
		\[f'(x) = 2(1 - x)\e^{-x}.\]
		
		\item Donner le tableau de variations de la fonction $f$ sur l'intervalle [0~;~1].
		
On considère la suite $\left(u_n\right)$ définie par $u_0 = 0,1$ et pour tout entier naturel $n$,
\[u_{n+1} = f\left(u_n\right).\]
	\end{enumerate}
\item 
	\begin{enumerate}
		\item Démontrer par récurrence que, pour tout $n$ entier naturel,
		\[0 \leqslant u_n < u_{n+1} \leqslant 1.\]
		
		\item En déduire que la suite $\left(u_n\right)$ est convergente.
	\end{enumerate}
\item Démontrer que la limite de la suite $\left(u_n\right)$ est $\ln (2)$.
\item 
	\begin{enumerate}
		\item Justifier que pour tout entier naturel $n,\: \ln (2) - u_n$ est positif.
		\item On souhaite écrire un script Python qui renvoie une valeur approchée de $\ln(2)$ par défaut à $10^{-4}$ près, ainsi que le nombre d'étapes pour y parvenir.
		
Recopier et compléter le script ci-dessous afin qu'il réponde au problème posé.

\begin{center}
\begin{tabular}{|l|}\hline
\textbf{def} seuil() :\\
\qquad n = 0\\
\qquad u = 0.1\\
\qquad \textbf{while} ln (2) $-$ u \ldots 0.0001 :\\
\qquad \qquad n=n+1\\
\qquad \qquad u= \ldots\\
\qquad \textbf{return} (u,~n)\\ \hline
\end{tabular}
\end{center}

		\item Donner la valeur de la variable $n$ renvoyée par la fonction seuil ().
	\end{enumerate}
\end{enumerate}

\medskip
